%% ============================================================
%%  Handout — Aula 14: Busca em árvore Monte Carlo II (AlphaZero)
%%  Complemento impresso dos slides aula14.tex
%%  Prof. Marcelo Keese Albertini — Universidade Federal de Uberlândia
%%  Compilar com XeLaTeX (2x): latexmk -xelatex handout14.tex
%% ============================================================
\documentclass[11pt]{article}

\usepackage{handoutAprendizadoRL}
\usepackage{babel}
\babelprovide[import, main]{brazilian}
\usepackage{multirow}

\emergencystretch=3em
\hyphenation{re-com-pen-sa re-com-pen-sas si-mu-la-ção si-mu-la-ções
  ar-re-pen-di-men-to re-tro-pro-pa-ga-ção pro-fun-di-da-de
  a-pren-di-za-gem re-gu-la-ri-za-da hi-per-pa-râ-me-tros
  im-ple-men-ta-ção mul-ti-pli-ca-ção com-bi-na-tó-ria
  Sze-pes-vá-ri Schritt-wie-ser Mas-te-ring Re-in-for-ce-ment}

\setaula{Aula 14 — Busca em árvore Monte Carlo II}
\setprof{Prof. Marcelo Keese Albertini}

\newcommand{\cpuct}{c_{\mathrm{puct}}}
\newcommand{\Nvis}{N}
\newcommand{\Wtot}{W}
\newcommand{\redeval}{v_{\theta}}
\newcommand{\redepol}{\bm{p}_{\theta}}
\newcommand{\raizpol}{\bm{\pi}}

\begin{document}

\capa{Busca em árvore Monte Carlo II: AlphaZero}
     {PUCT \textbullet\ auto-jogo \textbullet\ rede política--valor \textbullet\
      AlphaTensor e AlphaDev \textbullet\ raciocínio metanível}
     {Prof. Marcelo Keese Albertini}
     {Adaptado de \emph{CS234: Reinforcement Learning}, Emma Brunskill, Stanford University — Lecture 14}

\begin{ideiabox}
\textbf{Como usar este handout.} Ele não resume os slides: complementa-os. O que
no projetor aparece como resultado, aqui aparece com a passagem inteira — a
derivação da forma fechada por trás do PUCT, a conta da densidade de informação
no auto-jogo, o traço numérico completo da busca. As perguntas em
\textsf{Para discutir em aula} são as que o professor conduz oralmente; os
exercícios ao final incluem itens que não estão no CS234.
\end{ideiabox}

%% ============================================================
\section{Decidir agora, num mundo grande demais}

\subsection{Duas maneiras de responder ``o que eu faço?''}

Até esta altura do curso, quase todo algoritmo produziu o mesmo tipo de objeto:
uma política $\pol : \gS \to \gA$ definida sobre \emph{todo} o espaço de estados.
Iteração de valor, $Q$-learning, DQN, REINFORCE, PPO — todos entregam uma função
que responde a qualquer estado que apareça. Chame isso de resposta
\emph{amortizada}: paga-se caro uma vez, na fase de treino, e cada consulta
posterior é barata.

Existe outra maneira. Dado o estado atual $s_t$, gastar computação \emph{agora}
para escolher bem \emph{uma} ação, e só ela. Depois, no próximo estado, repetir do
zero. Chame isso de resposta \emph{sob demanda}: nada é pago antecipadamente, mas
cada decisão custa.

A escolha entre as duas não é estética. Quando $\abs{\gS}$ é astronômico — e
$10^{170}$ posições legais no Go é astronômico em qualquer escala razoável — a
primeira estratégia é literalmente impossível: não há memória no universo para
tabular $\pol$, nem dados para treinar uma aproximação uniformemente boa. Sobra a
segunda. E aí se descobre algo mais interessante: as duas não competem. AlphaZero
usa a busca local para \emph{fabricar dados} com que treina uma política global, e
usa a política global para \emph{acelerar} a busca local. O sistema inteiro é esse
laço.

\subsection{Busca Monte Carlo simples, e por que ela não basta}

O algoritmo mais direto: dado um simulador $\mathcal{M}$ e uma política de
simulação $\pol$, para cada ação $a$ simule $K$ episódios a partir de $s_t$ e
estime
\[
  \widehat{Q}(s_t,a) \;=\; \frac{1}{K}\sum_{k=1}^{K} G_t^{k}
  \;\xrightarrow[\;K\to\infty\;]{}\; Q^{\pol}(s_t,a),
  \qquad
  a_t = \argmax_{a} \widehat{Q}(s_t,a).
\]

Vale a pena ser explícito sobre o que isso é e o que não é. O estimador converge
para $Q^{\pol}$, não para $Q^{*}$. Escolher $\argmax_a Q^{\pol}(s_t,a)$ é
exatamente um passo de \emph{melhoria de política} sobre $\pol$. Pelo teorema da
melhoria de política, a política resultante é pelo menos tão boa quanto $\pol$ —
mas apenas um passo à frente dela.

\begin{alertabox}
Se $\pol$ for a política aleatória uniforme, a decisão será ``a melhor jogada
supondo que, a partir do lance seguinte, ambos os jogadores joguem ao acaso''. Em
Go isso é quase inútil: a maioria das posições tem valor semelhante sob jogo
aleatório. É por essa razão que a MCTS existe — para fazer \emph{mais} de um
passo de melhoria.
\end{alertabox}

\subsection{A árvore expectimax e a explosão combinatória}

Com um modelo, dá para calcular $Q^{*}(s_t,a)$ exatamente, expandindo a árvore
completa a partir de $s_t$ e retropropagando
\[
  Q^{*}(s,a) \;=\; \rec(s,a) \;+\; \gamma \sum_{s'} \tran{s'}{s,a}\,\max_{a'} Q^{*}(s',a').
\]
A observação útil é que \emph{não é preciso resolver o MDP inteiro}: basta o
sub-MDP que começa em $s_t$ e vai até o horizonte. A observação fatal é que esse
sub-MDP tem tamanho $(\abs{\gS}\abs{\gA})^{H}$.

\begin{center}
\begin{tabular}{lrrr}
\toprule
 & \textbf{ramificação} & \textbf{duração} & \textbf{folhas da árvore completa} \\
\midrule
Xadrez & $\approx 35$  & $\approx 80$  & $\approx 10^{123{,}5}$ \\
Go     & $\approx 250$ & $\approx 150$ & $\approx 10^{359{,}7}$ \\
\midrule
\multicolumn{3}{l}{\itshape átomos no universo observável} & $\approx 10^{80}$ \\
\bottomrule
\end{tabular}
\end{center}

\subsection{MCTS: crescer a árvore só onde importa}

A resposta da busca em árvore Monte Carlo é não construir a árvore toda. Constrói-se
uma árvore \emph{assimétrica}, aprofundada apenas nas regiões que as simulações
anteriores indicaram valer a pena. Cada simulação tem quatro fases:

\begin{enumerate}
  \item \textbf{Seleção.} A partir da raiz, descer pela árvore existente
        escolhendo, em cada nó, a ação que maximiza um critério que combina valor
        estimado e incerteza.
  \item \textbf{Expansão.} Ao chegar a um nó que ainda não está na árvore,
        acrescentá-lo.
  \item \textbf{Avaliação.} Obter um valor $v$ para esse nó novo — por rollout
        (MCTS clássica) ou por uma rede (AlphaZero).
  \item \textbf{Retropropagação.} Subir pelo caminho percorrido atualizando, em
        cada aresta $(s,a)$: $\Nvis \mathrel{+}= 1$, $\Wtot \mathrel{+}= v$,
        $Q = \Wtot/\Nvis$.
\end{enumerate}

\begin{ideiabox}
Todo o projeto de um algoritmo de MCTS está em \textbf{duas decisões}: o critério
da fase 1 e a origem do $v$ da fase 3. Tudo o mais é contabilidade. Vale a pena
carregar esse par como lente ao ler qualquer variante da literatura.
\end{ideiabox}

\subsection{UCT, e um detalhe a guardar}

O UCT (Kocsis e Szepesvári, 2006) resolve a fase 1 tratando cada nó como um
bandit de múltiplos braços:
\[
  a_{ik} \;=\; \argmax_{a}\ \Biggl[\;
  \frac{1}{\Nvis(i,a)}\sum_{k=1}^{\Nvis(i,a)} G_k(i,a)
  \;+\; c\sqrt{\frac{\log \Nvis(i)}{\Nvis(i,a)}}\;\Biggr].
\]

Duas consequências merecem destaque, porque voltarão na Seção~\ref{sec:critica}.

Primeiro, \textbf{a política de simulação muda a cada episódio}: a árvore não é
percorrida por uma política fixa, mas por uma que se adapta ao histórico. Segundo,
e como corolário direto, \textbf{a média $\overline{G}$ retropropagada é uma média
sobre uma política não estacionária}. Ela não estima $Q^{*}$ nem $Q^{\pol}$ para
nenhum $\pol$ fixo; estima a média histórica de uma sequência de políticas cada
vez melhores. Converge para o valor ótimo, sim, mas por baixo e devagar — e o
atraso de cada nível se compõe ao subir pela árvore.

%% ============================================================
\section{AlphaZero: a busca guiada por uma rede}

\subsection{Uma rede, duas cabeças}

Uma única rede residual $f_{\theta}$ recebe a posição $s$ e devolve o par
\[
  f_{\theta}(s) \;=\; \bigl(\redepol(s),\ \redeval(s)\bigr),
  \qquad \redepol(s) \in \Delta(\gA), \qquad \redeval(s) \in [-1,1].
\]
A cabeça de política é usada como \emph{prior} na fase de seleção; a cabeça de
valor substitui a rollout na fase de avaliação. O corpo é compartilhado porque as
duas tarefas exigem as mesmas representações: reconhecer grupos vivos e mortos,
influência, forma.

\begin{teobox}{Uma passagem direta contra cem rollouts}
No AlphaGo de 2016, uma única passagem de $\redeval$ atingiu erro quadrático médio
comparável à média de 100 rollouts com a política de RL — a um custo ordens de
grandeza menor. Esse número explica a trajetória posterior: a otimização bayesiana
dos hiperparâmetros do AlphaGo deslocou progressivamente o peso da mistura em
direção à rede de valor a cada ciclo de projeto, até a rollout ser removida por
completo no AlphaGo Zero (Chen et al., 2018).
\end{teobox}

\subsection{PUCT: anatomia da fórmula}

Em cada nó, mantêm-se por ação: $\Nvis(s,a)$, $\Wtot(s,a)$, $Q(s,a) = \Wtot/\Nvis$
e o prior $P(s,a) = \redepol(s)_a$. A ação a simular é

\begin{teobox}{Regra de seleção do AlphaZero}
\[
  a_t \;=\; \argmax_{a}\ \bigl[\, Q(s,a) + U(s,a) \,\bigr],
  \qquad
  U(s,a) \;=\; \cpuct\, P(s,a)\, \frac{\sqrt{\sum_b \Nvis(s,b)}}{1+\Nvis(s,a)} .
\]
\end{teobox}

Leia termo a termo. O numerador $\sqrt{\Nvis(s)}$ faz o bônus \emph{crescer} com a
quantidade total de trabalho já feito no nó — quanto mais o nó foi visitado, mais
se justifica reexaminar alternativas. O denominador $1+\Nvis(s,a)$ faz o bônus
\emph{cair} com o trabalho já feito naquela ação específica. O fator $P(s,a)$
pondera tudo pela opinião da rede. E o $+1$ no denominador garante que uma ação
nunca visitada tenha bônus \emph{finito}, igual a $\cpuct P(s,a)\sqrt{\Nvis(s)}$ —
diferente do UCT, em que ações não visitadas têm bônus infinito e são todas
tentadas obrigatoriamente.

\begin{center}
\begin{tabular}{lll}
\toprule
 & \textbf{UCT} & \textbf{PUCT} \\
\midrule
Bônus & $c\sqrt{\log \Nvis(s)/\Nvis(s,a)}$ & $\cpuct P(s,a)\sqrt{\Nvis(s)}/(1+\Nvis(s,a))$ \\
Decaimento em $\Nvis(s,a)$ & $\Nvis(s,a)^{-1/2}$ & $\Nvis(s,a)^{-1}$ \\
Crescimento em $\Nvis(s)$ & $\sqrt{\log \Nvis(s)}$ & $\sqrt{\Nvis(s)}$ \\
Ação nunca visitada & bônus infinito & bônus finito, $\propto P(s,a)$ \\
Justificativa & desigualdade de Hoeffding & heurística \\
Avaliação da folha & rollout até o fim & $\redeval$ \\
\bottomrule
\end{tabular}
\end{center}

\begin{alertabox}
\textbf{O PUCT não é uma cota superior de confiança.} Não há desigualdade de
concentração por trás dele, e a consequência é operacional: \emph{ações com prior
baixo podem nunca ser visitadas}. Em Go, com $\approx 250$ jogadas legais e
$1\,600$ simulações, é isso que torna a busca viável — e é isso que exige o ruído
de Dirichlet da Seção~\ref{sec:dir}.
\end{alertabox}

\subsection{O PUCT em câmera lenta}
\label{sec:traco}

Considere uma raiz com três jogadas em que a rede erra: dá prior alto para $a_1$,
que é a pior das duas jogadas razoáveis. Priors $P = (0{,}60;\ 0{,}30;\ 0{,}10)$;
valores das folhas $\redeval = (0{,}10;\ 0{,}55;\ -0{,}40)$; $\cpuct = 1{,}5$.
Na primeira simulação, com $\Nvis(s)=0$, adota-se a convenção
$\sqrt{\Nvis(s)} \to 1$.

\begin{center}
\small
\begin{tabular}{c ccc c ccc c ccc c}
\toprule
& \multicolumn{3}{c}{$Q(s,a)$} && \multicolumn{3}{c}{$U(s,a)$} && \multicolumn{3}{c}{$Q+U$} & \\
\cmidrule(lr){2-4}\cmidrule(lr){6-8}\cmidrule(lr){10-12}
$k$ & $a_1$ & $a_2$ & $a_3$ && $a_1$ & $a_2$ & $a_3$ && $a_1$ & $a_2$ & $a_3$ & escolhe \\
\midrule
 1 & 0,00 & 0,00 & 0,00 && 0,900 & 0,450 & 0,150 && \textbf{0,900} & 0,450 & 0,150 & $a_1$ \\
 2 & 0,10 & 0,00 & 0,00 && 0,450 & 0,450 & 0,150 && \textbf{0,550} & 0,450 & 0,150 & $a_1$ \\
 3 & 0,10 & 0,00 & 0,00 && 0,424 & 0,636 & 0,212 && 0,524 & \textbf{0,636} & 0,212 & $a_2$ \\
 4 & 0,10 & 0,55 & 0,00 && 0,520 & 0,390 & 0,260 && 0,620 & \textbf{0,940} & 0,260 & $a_2$ \\
 5 & 0,10 & 0,55 & 0,00 && 0,600 & 0,300 & 0,300 && 0,700 & \textbf{0,850} & 0,300 & $a_2$ \\
 6 & 0,10 & 0,55 & 0,00 && 0,671 & 0,252 & 0,335 && 0,771 & \textbf{0,802} & 0,335 & $a_2$ \\
 7 & 0,10 & 0,55 & 0,00 && 0,735 & 0,220 & 0,367 && \textbf{0,835} & 0,770 & 0,367 & $a_1$ \\
 8 & 0,10 & 0,55 & 0,00 && 0,595 & 0,238 & 0,397 && 0,695 & \textbf{0,788} & 0,397 & $a_2$ \\
 9 & 0,10 & 0,55 & 0,00 && 0,636 & 0,212 & 0,424 && 0,736 & \textbf{0,762} & 0,424 & $a_2$ \\
10 & 0,10 & 0,55 & 0,00 && 0,675 & 0,193 & 0,450 && \textbf{0,775} & 0,743 & 0,450 & $a_1$ \\
\bottomrule
\end{tabular}
\end{center}

Quatro leituras, na ordem em que aparecem.

\textbf{Simulações 1--2: a rede manda sozinha.} Nada foi avaliado, $Q \equiv 0$, e
o ranking é o do prior. Esta é a fase em que a qualidade de $\redepol$ é tudo.

\textbf{Simulação 3: a evidência entra.} Assim que $a_1$ devolve $0{,}10$ —
abaixo do que o prior insinuava — o somatório muda de líder. Note que bastou
\emph{uma} avaliação para virar o jogo; é o termo $U$ decaindo como
$\Nvis(s,a)^{-1}$ que permite essa rapidez.

\textbf{Simulação 7: a busca volta a conferir $a_1$.} O denominador
$1+\Nvis(s,a_2)$ cresceu o bastante para que $a_1$ recupere a liderança
momentaneamente. Isto não é desperdício: é o mecanismo que impede a busca de se
fechar em torno de um erro precoce.

\textbf{A ação $a_3$ fica fora até $k=23$.} Prior baixo mais orçamento pequeno é
igual a jogada nunca examinada. Com 16 simulações, $\Nvis = (5,\ 11,\ 0)$.

\subsection{Quando o prior está muito errado}

O traço acima é benigno: o prior erra, mas não muito. O caso preocupante é o prior
confiante e errado. Mantendo $\redeval$ e trocando os priors para
$P = (0{,}95;\ 0{,}04;\ 0{,}01)$:

\begin{center}
\begin{tabular}{rccc l}
\toprule
simulações & $\Nvis(a_1)$ & $\Nvis(a_2)$ & $\Nvis(a_3)$ & recomendação da raiz \\
\midrule
 50 & 33 &  17 & 0 & \textbf{$a_1$ — errado} \\
200 & 44 & 156 & 0 & $a_2$ \\
800 & 89 & 711 & 0 & $a_2$ \\
\bottomrule
\end{tabular}
\end{center}

A busca corrige a rede, mas o preço é orçamento — e por um bom tempo ela
\emph{confirma} o erro em vez de corrigi-lo. Esta é a formulação quantitativa do
princípio geral: rede melhor significa menos simulações para a mesma decisão. É
também a razão pela qual a curva de ``força em função do número de simulações''
tem um joelho: até certo ponto, mais busca compra correção do prior; depois disso,
compra pouco.

\subsection{$\cpuct$: quanto duvidar da própria rede}

Mesmo exemplo, orçamento fixo de 400 simulações, variando só $\cpuct$:

\begin{center}
\begin{tabular}{ccccl}
\toprule
$\cpuct$ & $\Nvis(a_1)$ & $\Nvis(a_2)$ & $\Nvis(a_3)$ & interpretação \\
\midrule
0,5 &  13 & 386 & 1 & confia na rede; converge rápido \\
1,0 &  25 & 373 & 2 & \\
1,5 &  37 & 360 & 3 & faixa típica na prática \\
3,0 &  71 & 324 & 5 & \\
5,0 & 107 & 284 & 9 & duvida da rede; gasta reexaminando o descartado \\
\bottomrule
\end{tabular}
\end{center}

\begin{ideiabox}
Não é um parâmetro cosmético. No AlphaGo, ajustar por otimização bayesiana os
hiperparâmetros da MCTS — $\cpuct$ entre eles — rendeu \textbf{94 e 103 pontos de
Elo} em dois ciclos de projeto antes da partida contra Lee Sedol. Cem pontos de
Elo correspondem a cerca de 64\% de vitórias contra a versão anterior. Ao ler
comparações entre algoritmos de MCTS na literatura, pergunte sempre quanto de
ajuste houve de cada lado.
\end{ideiabox}

\subsection{Ruído de Dirichlet}
\label{sec:dir}

Se o prior é o que abre a porta, e a porta de $a_3$ nunca abre, o sistema não tem
como descobrir que a rede está errada sobre $a_3$. Em auto-jogo — e apenas em
auto-jogo — o AlphaZero perturba o prior na raiz:
\[
  P(s_0,a) \;=\; (1-\varepsilon)\,\redepol(s_0)_a \;+\; \varepsilon\,\eta_a,
  \qquad \eta \sim \mathrm{Dir}(\alpha), \qquad \varepsilon = 0{,}25 .
\]

O parâmetro $\alpha$ é calibrado inversamente ao número típico de jogadas legais:
$0{,}30$ para xadrez ($\approx 35$ jogadas), $0{,}15$ para shogi ($\approx 90$) e
$0{,}03$ para Go ($\approx 250$).

Vale entender por que $\alpha$ pequeno. Uma Dirichlet com $\alpha \ll 1$ produz
vetores \emph{esparsos}: quase toda a massa cai em uma ou duas coordenadas
sorteadas. O efeito é que, em cada partida de auto-jogo, um punhado de jogadas
recebe um empurrão \emph{grande}, em vez de todas receberem um empurrão
minúsculo. Isso gera exploração profunda — a jogada sorteada é realmente
investigada — em vez de uma diluição uniforme que não muda nada.

\begin{ideiabox}
O ruído fica só na raiz, e isso é essencial. A busca precisa \emph{poder} tentar a
jogada estranha, mas continua livre para descartá-la assim que $Q$ mostrar que é
ruim. É exploração que a evidência pode vetar — diferente de $\epsilon$-guloso,
que força a execução da ação aleatória. Em avaliação e partidas oficiais, o ruído
é removido: ele é um mecanismo de \emph{coleta de dados}, não de jogo.
\end{ideiabox}

\subsection{A política da raiz e a temperatura}

Terminadas as $K$ simulações, a política exportada é
\[
  \raizpol(a \mid s_0) \;=\; \frac{\Nvis(s_0,a)^{1/\tau}}{\sum_b \Nvis(s_0,b)^{1/\tau}} .
\]
Com $\Nvis = (5,11,0)$ do traço da Seção~\ref{sec:traco}:

\begin{center}
\begin{tabular}{ccccc}
\toprule
$\tau$ & $\raizpol(a_1)$ & $\raizpol(a_2)$ & $\raizpol(a_3)$ & $H[\raizpol]$ (bits) \\
\midrule
1,00 & 0,3125 & 0,6875 & 0 & 0,896 \\
0,50 & 0,1712 & 0,8288 & 0 & 0,661 \\
0,25 & 0,0409 & 0,9591 & 0 & 0,247 \\
0,10 & 0,0004 & 0,9996 & 0 & 0,005 \\
\bottomrule
\end{tabular}
\end{center}

No auto-jogo usa-se $\tau = 1$ nas 30 primeiras jogadas — para gerar diversidade
de posições no conjunto de treino — e $\tau \to 0$ depois, para jogar a partida
com seriedade.

Por que exportar contagens de visita em vez de $\argmax_a Q(s_0,a)$, que parece a
estimativa mais direta? Porque $\Nvis$ é um resumo mais robusto. Uma ação visitada
três vezes pode ter $Q$ alto por sorte; $Q$ sozinho não carrega informação sobre
quanta evidência o sustenta. As contagens já incorporam a confiança: só recebe
muitas visitas o que sobreviveu à comparação repetida com as alternativas ao longo
de toda a busca.

\subsection{O que o PUCT está de fato otimizando}

As contagens de visita não são um subproduto arbitrário da regra de seleção.
Grill et al.\ (ICML 2020) mostraram que a distribuição de visitas normalizada
resolve aproximadamente um problema de otimização de política regularizada por
divergência de Kullback--Leibler.

\begin{provabox}
\textbf{Problema.} Com $\bm{q}$ os valores estimados e $\redepol$ o prior,
considere
\[
  \bar{\raizpol} \;=\; \argmax_{y \in \Delta(\gA)}
  \Bigl[\, \bm{q}^{\top} y \;-\; \lambda\,\mathrm{KL}(\redepol \,\|\, y) \Bigr],
  \qquad \lambda > 0 .
\]

\textbf{Derivação.} Escrevendo
$\mathrm{KL}(\redepol\|y) = \sum_a p_a \log p_a - \sum_a p_a \log y_a$, o primeiro
somatório não depende de $y$ e pode ser descartado. O lagrangiano com o vínculo
$\sum_a y_a = 1$ é
\[
  \mathcal{L}(y,\alpha) \;=\; \sum_a q_a y_a \;+\; \lambda \sum_a p_a \log y_a
  \;-\; \alpha\Bigl(\sum_a y_a - 1\Bigr).
\]
Derivando em $y_a$ e igualando a zero:
\[
  q_a + \frac{\lambda p_a}{y_a} - \alpha = 0
  \qquad\Longrightarrow\qquad
  \boxed{\ y_a = \frac{\lambda\, p_a}{\alpha - q_a}\ }
\]
com $\alpha$ determinado pela normalização
$\sum_a \lambda p_a/(\alpha - q_a) = 1$. Como o logaritmo força $y_a > 0$, é
necessário $\alpha > \max_a q_a$; o lado esquerdo é estritamente decrescente em
$\alpha$ nesse intervalo e vai de $+\infty$ a $0$, logo a raiz existe e é única.

\textbf{Casos-limite.} Quando $\lambda \to \infty$, a normalização exige
$\alpha \to \infty$ e $y_a \to p_a$: recupera-se o prior. Quando
$\lambda \to 0^{+}$, tem-se $\alpha \to (\max_a q_a)^{+}$, e a coordenada que
atinge o máximo domina: $y \to$ indicadora do $\argmax$ de $\bm{q}$.
\end{provabox}

A conexão com o PUCT é a identificação
$\lambda_{\Nvis} = \cpuct\sqrt{\Nvis}/(\Nvis + \abs{\gA})$ e
$\bar{\raizpol}_a \approx \Nvis(s,a)/(\Nvis + \abs{\gA})$. A tabela abaixo executa
o PUCT do exemplo da Seção~\ref{sec:traco} e compara com a forma fechada:

\begin{center}
\small
\begin{tabular}{r ccc c ccc c c}
\toprule
& \multicolumn{3}{c}{visitas empíricas} && \multicolumn{3}{c}{forma fechada} && \\
\cmidrule(lr){2-4}\cmidrule(lr){6-8}
$\Nvis$ & $a_1$ & $a_2$ & $a_3$ && $a_1$ & $a_2$ & $a_3$ && $\lambda_{\Nvis}$ \\
\midrule
   64 & 0,1940 & 0,7463 & 0,0149 && 0,2069 & 0,7755 & 0,0176 && 0,1791 \\
  256 & 0,1120 & 0,8687 & 0,0077 && 0,1154 & 0,8752 & 0,0094 && 0,0927 \\
 1600 & 0,0480 & 0,9464 & 0,0037 && 0,0486 & 0,9475 & 0,0039 && 0,0374 \\
 6400 & 0,0245 & 0,9731 & 0,0019 && 0,0247 & 0,9734 & 0,0020 && 0,0187 \\
\bottomrule
\end{tabular}
\end{center}

O acordo é bom já em $\Nvis = 256$ e excelente daí em diante. A leitura conceitual:
como $\lambda_{\Nvis} \sim \Nvis^{-1/2} \to 0$, a regularização enfraquece com o
orçamento. Pouco orçamento significa ``fique perto do que a rede acha''; muito
orçamento significa ``vá para o $\argmax$ dos valores''. É a mesma estrutura de
proximidade regularizada do TRPO e do PPO — só que aqui ela emerge de um
planejador, sem ter sido escrita em lugar nenhum.

\begin{discbox}
O PUCT foi projetado como heurística e só treze anos depois foi reinterpretado
como otimização regularizada. Isso é comum em aprendizagem de máquina. Que outros
exemplos do curso seguem o mesmo padrão — heurística primeiro, justificativa
depois? A justificativa \emph{a posteriori} tem valor prático, ou é apenas
arrumação estética?
\end{discbox}

%% ============================================================
\section{Auto-jogo: o laço que fabrica os próprios dados}

\subsection{A estrutura do laço}

O sistema completo é um ciclo de quatro etapas, sem nenhum humano dentro:
a busca guiada por $f_\theta$ produz $\raizpol$; a partida é jogada amostrando de
$\raizpol$; cada jogada vira um exemplo $(s_t, \raizpol_t, z)$; a rede é treinada
nesses exemplos; e a rede nova volta a guiar a busca.

O único insumo externo são as \emph{regras do jogo}, usadas pelo simulador para
gerar sucessores e detectar o fim da partida.

\begin{center}
\begin{tabular}{ll}
\toprule
Partidas por iteração & 25\,000 \\
Simulações por jogada & 1\,600 ($\approx 0{,}4$\,s por busca) \\
Temperatura & $\tau = 1$ nas 30 primeiras jogadas, depois $\tau \to 0$ \\
Ruído na raiz & $\varepsilon = 0{,}25$, $\eta \sim \mathrm{Dir}(0{,}03)$ no Go \\
Desistência & automática, mantendo falsos positivos abaixo de 5\% \\
Rodada de 3 dias & 4,9 milhões de partidas, rede de 20 blocos \\
Rodada de 40 dias & 29 milhões de partidas, rede de 40 blocos, Elo 5\,185 \\
\bottomrule
\end{tabular}
\end{center}

\subsection{A perda, e a assimetria entre os dois alvos}

\begin{teobox}{Função de perda}
\[
  \ell(\theta) \;=\;
  \underbrace{\bigl(z - \redeval(s)\bigr)^{2}}_{\text{regressão do valor}}
  \;-\; \underbrace{\raizpol^{\top}\log \redepol(s)}_{\text{entropia cruzada}}
  \;+\; \underbrace{c\,\norm{\theta}^{2}}_{\text{regularização}}
\]
\end{teobox}

Os dois alvos têm naturezas muito diferentes, e confundi-los é fonte comum de
mal-entendidos.

O \textbf{alvo de política} $\raizpol_t$ é \emph{denso e rico}. Uma partida de 200
jogadas produz 200 alvos, e cada um é uma distribuição inteira sobre as jogadas
legais — não um rótulo duro. É o resultado do passo de melhoria sendo destilado na
rede.

O \textbf{alvo de valor} $z$ é \emph{escasso e correlacionado}: um único bit por
partida, replicado em todas as posições dela. Duas posições da mesma partida
compartilham exatamente o mesmo rótulo, ainda que uma seja do lance 10 e a outra
do lance 190.

\begin{alertabox}
Diz-se com frequência que ``Go é um problema de recompensa esparsa''. É meia
verdade. A \emph{recompensa} é terminal, mas o \emph{alvo de treino da política} é
denso, porque a busca fabrica supervisão em cada posição. É um dos motivos pelos
quais o auto-jogo com MCTS aprende, enquanto gradiente de política puro com
recompensa terminal patinaria: a MCTS converte um sinal terminal em supervisão por
jogada.
\end{alertabox}

\subsection{Currículo automático e densidade de informação}

O adversário do AlphaZero é uma cópia de si mesmo. A consequência é que a taxa de
vitória fica em torno de 50\% durante todo o treino, e isso é mais importante do
que parece.

Trate o resultado $z$ como uma variável de Bernoulli com parâmetro $p$ igual à
probabilidade de vitória. A informação que uma partida carrega sobre os parâmetros
é máxima quando a entropia de $z$ é máxima:

\begin{center}
\begin{tabular}{ccc}
\toprule
$p$ (prob.\ de vitória) & $H(z)$ em bits & partidas por bit de informação \\
\midrule
0,50 & 1,000 & 1,0 \\
0,75 & 0,811 & 1,2 \\
0,90 & 0,469 & 2,1 \\
0,99 & 0,081 & 12,4 \\
\bottomrule
\end{tabular}
\end{center}

Contra um adversário fixo e muito mais forte, todas as partidas terminariam em
derrota: $z \equiv -1$, entropia zero, gradiente do termo de valor praticamente
nulo. Contra um muito mais fraco, o mesmo com o sinal trocado. Um adversário
sempre à altura é a definição operacional de um \emph{currículo}: a dificuldade
sobe exatamente na velocidade da competência.

\begin{ideiabox}
Esta é uma resposta parcial — e desconfortável — à pergunta ``por que RL funciona
tão bem em jogos e sofre tanto no mundo real''. Jogos de dois jogadores vêm com um
gerador de currículo embutido, de graça. Robótica, medicina e diálogo não vêm.
Onde não há auto-jogo, o currículo tem de ser construído à mão, e isso é
engenharia difícil.
\end{ideiabox}

\begin{discbox}
Existem domínios não competitivos em que se pode fabricar um análogo do auto-jogo?
Pense em geração de problemas cuja dificuldade acompanha a competência atual do
agente — provas matemáticas, casos clínicos sintéticos, cenários de tráfego. O que
faz o papel do ``adversário à altura'' nesses casos, e o que se perde em relação
ao jogo de dois jogadores?
\end{discbox}

%% ============================================================
\section{O que a evidência mostra}

O artigo do AlphaGo Zero é incomum por trazer ablações limpas. Três perguntas
foram isoladas: quanto do avanço vem da arquitetura, quanto vem da busca, e quanto
se perde ao dispensar dados humanos.

\subsection{Arquitetura}

Quatro redes foram treinadas no mesmo conjunto fixo de partidas de auto-jogo,
combinando política e valor \emph{separadas} (sep) ou \emph{conjuntas} (dual), e
arquitetura \emph{convolucional} (conv) ou \emph{residual} (res). Todas foram
avaliadas com a mesma busca, 5\,s por jogada.

\begin{center}
\begin{tabular}{lll}
\toprule
Configuração & Ganho em Elo & Corresponde a \\
\midrule
sep--conv  & base       & arquitetura do AlphaGo Lee \\
dual--conv & $+600$     & \\
sep--res   & $+600$     & \\
dual--res  & $+1\,200$  & arquitetura do AlphaGo Zero \\
\bottomrule
\end{tabular}
\end{center}

Cada uma das duas mudanças vale, sozinha, mais de 600 Elo; juntas, mais de 1\,200 —
a distância entre um forte amador e um profissional. É um resultado desconfortável
e honesto: parte substancial do avanço rotineiramente atribuído ``ao algoritmo''
era engenharia de rede. Curiosamente, a cabeça conjunta \emph{piora} ligeiramente
a previsão de jogadas humanas; em compensação reduz o erro de valor e fortalece o
jogo. Prever bem o humano e jogar bem não são o mesmo objetivo.

\subsection{Busca}

A mesma rede treinada, jogando sem nenhum lookahead — apenas
$\argmax_a \redepol(s)_a$ — contra o sistema completo:

\begin{center}
\begin{tabular}{lrl}
\toprule
Jogador & Elo & \\
\midrule
Rede crua do AlphaGo Zero (sem busca) & 3\,055 & \\
AlphaGo Fan   & 3\,144 & venceu Fan Hui \\
AlphaGo Lee   & 3\,739 & venceu Lee Sedol \\
AlphaGo Master & 4\,858 & 60--0 contra profissionais on-line \\
AlphaGo Zero  & 5\,185 & 100--0 contra o AlphaGo Lee \\
\bottomrule
\end{tabular}
\end{center}

A diferença entre a rede crua e o sistema completo é de $2\,130$ pontos. Na escala
Elo, a probabilidade de vitória do mais forte é
$p = 1/(1 + 10^{-\Delta/400})$, o que dá:

\begin{center}
\begin{tabular}{rrr}
\toprule
$\Delta$ Elo & $p$(vitória) & derrota a cada \\
\midrule
  200 & 0,7597 & 4 partidas \\
  600 & 0,9693 & 33 partidas \\
1\,200 & 0,9990 & 1\,001 partidas \\
2\,130 & 0,999995 & 211\,350 partidas \\
\bottomrule
\end{tabular}
\end{center}

\begin{teobox}{Busca é computação em tempo de decisão}
A mesma rede, com e sem MCTS, difere em mais de 2\,000 pontos de Elo. Não houve
treino adicional, nem parâmetro novo: apenas o direito de pensar por 0,4 segundo
antes de responder. É a mesma alavanca hoje explorada em modelos de linguagem sob
o nome de \emph{test-time compute}.
\end{teobox}

\subsection{Dados humanos}

A comparação direta rodou o mesmo laço partindo de uma rede inicializada por
aprendizado supervisionado sobre partidas humanas (conjunto KGS) e partindo de
pesos aleatórios. A versão supervisionada aprende mais rápido \emph{no início} e
prevê jogadas humanas com mais acurácia \emph{o tempo todo}; a versão
\emph{tabula rasa} a ultrapassa em poucas horas e termina mais forte.

\begin{alertabox}
Cuidado com a generalização apressada. Aqui existem duas condições raras: um
\emph{simulador perfeito} e uma \emph{recompensa verificável} ao final. Onde
faltam essas duas coisas — diálogo, medicina, direção autônoma — a conclusão
``dados humanos são dispensáveis'' não se transfere. O que se transfere é a
observação mais fina: imitação carrega o teto do imitado, e o auto-jogo não tem
esse teto \emph{quando existe um juiz confiável}.
\end{alertabox}

\subsection{Um algoritmo, três jogos}

O AlphaZero removeu o que restava de específico do Go — simetrias do tabuleiro,
torneio de avaliação entre gerações — e rodou o mesmo código em xadrez, shogi e Go.

\begin{center}
\begin{tabular}{lll}
\toprule
 & \textbf{AlphaZero} & \textbf{Stockfish 8} \\
\midrule
Busca & MCTS guiada por rede & alfa-beta com poda \\
Posições por segundo & $\approx 6\times10^{4}$ & $\approx 6\times10^{7}$ \\
Conhecimento embutido & regras do xadrez & aberturas, avaliação, finais \\
\bottomrule
\end{tabular}
\end{center}

No confronto de mil partidas publicado na \emph{Science} em 2018: 155 vitórias,
6 derrotas, 839 empates. O treino levou 700\,000 passos de gradiente, e o
AlphaZero ultrapassou o Stockfish em cerca de quatro horas.

Mil vezes menos posições examinadas, e ainda assim vitória. A rede não acelera a
busca; ela escolhe muito melhor o que examinar. Buscar menos e melhor bateu buscar
muito e cego.

\begin{alertabox}
As condições da partida foram legitimamente contestadas: versão do Stockfish,
hardware assimétrico, controle de tempo, ausência de livro de aberturas. Em
partidas com desvantagem de tempo, o AlphaZero se manteve superior até a proporção
de 10:1. A comparação é informativa sobre \emph{paradigmas de busca}; não é uma
medição limpa de qual programa é mais forte.
\end{alertabox}

%% ============================================================
\section{Além do tabuleiro}

A receita, escrita sem menção a jogos: \emph{se um problema pode ser formulado
como sequência de escolhas com um simulador que verifica o resultado, então uma
rede que propõe candidatos somada a uma busca que testa candidatos pode superar a
heurística humana, mesmo em espaços astronômicos.} A única peça que muda entre
domínios é o simulador.

\subsection{AlphaTensor}

Multiplicar duas matrizes $n \times n$ equivale a decompor um tensor
$\mathcal{T}_n$ em soma de tríades de posto 1; o número de parcelas é o número de
multiplicações escalares. Cada ação do jogo é uma tríade, e o episódio termina
quando o tensor residual zera.

\begin{center}
\begin{tabular}{lccl}
\toprule
Problema & melhor conhecido & AlphaTensor & observação \\
\midrule
$2\times2$ (Strassen, 1969) & 7 & 7 & redescoberto do zero \\
$3\times3$ (Laderman, 1976) & 23 & 23 & redescoberto do zero \\
$4\times4$ sobre $\mathbb{Z}_2$ & 49 & 47 & primeiro avanço desde 1969 \\
$4\times5$ por $5\times5$ (reais) & 80 & 76 & \\
$5\times5$ (reais) & 98 & 96 & \\
\bottomrule
\end{tabular}
\end{center}

Duas ressalvas que a cobertura na imprensa perdeu. O resultado do $4\times4$ vale
em aritmética modular ($\mathbb{Z}_2$), não sobre os reais. E, meses depois,
Kauers e Moosbauer baixaram o caso $5\times5$ sobre GF(2) a 95 multiplicações
usando busca combinatória clássica — a última palavra ainda não é da máquina.

Um detalhe metodológico interessante: quando a recompensa foi modificada para
penalizar o \emph{tempo medido} numa GPU V100 específica, o sistema encontrou
algoritmos 10--20\% mais rápidos \emph{naquele} hardware. Otimalidade aqui é
relativa à máquina, não uma propriedade matemática.

\subsection{AlphaDev}

No \emph{AssemblyGame}, o estado é o programa em assembly construído até agora, a
ação é acrescentar uma instrução, e a recompensa combina correção (testada contra
entradas) com latência.

O alvo foram \texttt{sort3}, \texttt{sort4} e \texttt{sort5} — rotinas chamadas
por dentro de qualquer ordenação maior. O sistema descobriu duas sequências novas,
batizadas de \emph{swap} e \emph{copy} do AlphaDev, que economizam uma instrução
cada vez que são aplicadas; para \texttt{VarSort4} encontrou um algoritmo com
estrutura diferente e 29 instruções a menos que o de referência humano. Os
resultados foram integrados à \texttt{libc++} do LLVM — a primeira mudança nessa
parte da biblioteca em mais de uma década, e a primeira originada de aprendizagem
por reforço.

Os ganhos: até 70\% em sequências de cinco elementos, $\approx 1{,}7\%$ em
sequências acima de $250\,000$ elementos, e $\approx 30\%$ em funções de
\emph{hash} na faixa de 9 a 16 bytes. A discrepância entre 70\% e 1,7\% não é
contradição: as rotinas otimizadas são as \emph{folhas} de algoritmos maiores.
Melhorar a folha melhora muito a folha e pouco o todo — mas a folha é chamada
trilhões de vezes por dia.

\subsection{MuZero}

AlphaZero exige um simulador perfeito. O MuZero remove essa exigência aprendendo
um modelo — mas apenas o suficiente para planejar. Três funções são treinadas
conjuntamente: representação $h$ (observação $\mapsto$ estado latente), dinâmica
$g$ ($(s^k,a) \mapsto (s^{k+1}, r^k)$, no espaço latente) e predição $f$
($s^k \mapsto (\bm{p}, v)$). A MCTS roda inteiramente no espaço latente.

\begin{ideiabox}
\textbf{Modelo suficiente para decidir, não para prever.} Nada obriga $s^k$ a
reconstruir a observação: o modelo é treinado apenas para que política, valor e
recompensa saiam certos. É a diferença entre um modelo \emph{gerativo} do mundo e
um modelo \emph{orientado ao valor} — e prever pixels é caro e, em boa parte,
irrelevante para a decisão.
\end{ideiabox}

%% ============================================================
\section{Um olhar crítico: o UCT é o algoritmo certo?}
\label{sec:critica}

\subsection{A pergunta incômoda}

O UCT trata cada nó como um bandit e usa uma cota superior de confiança para
escolher o que simular. Por que isso é estranho?

Em um bandit \emph{real}, cada braço puxado custa: a recompensa perdida é perda de
verdade, e minimizar arrependimento cumulativo é exatamente o objetivo. Dentro da
simulação, nada disso vale. Puxar mil vezes um braço ruim custa apenas tempo de
CPU; a ``recompensa'' simulada não entra no bolso de ninguém. O objetivo correto
não é acumular recompensa simulada — é \emph{tomar a melhor decisão na raiz ao
final do orçamento}.

\begin{center}
\begin{tabular}{ll}
\toprule
\textbf{Arrependimento cumulativo} & \textbf{Arrependimento simples} \\
$R_n = \sum_{t=1}^{n}(\mu^{*} - \mu_{A_t})$ &
$r_n = \mu^{*} - \mu_{\widehat{A}_n}$ \\[0.3em]
o que se pagou pelo caminho & a qualidade da recomendação final \\
objetivo quando as ações são reais & objetivo quando as ações são simuladas \\
\bottomrule
\end{tabular}
\end{center}

\begin{teobox}{Bubeck, Munos e Stoltz (2009): existe um conflito}
Quanto mais um algoritmo minimiza o arrependimento \emph{cumulativo}, pior fica
seu arrependimento \emph{simples} no pior caso. Alocação uniforme atinge
$r_n = O(e^{-cn})$. Já para qualquer estratégia com arrependimento cumulativo
limitado por $C n^{a}$ uniformemente sobre as instâncias, existe uma instância em
que $r_n \ge \tfrac{1}{2K}\exp(-C' n^{a})$. Para o UCB, com $R_n = O(\log n)$, o
expoente efetivo é logarítmico: o arrependimento simples decai apenas
\emph{polinomialmente} em $n$.
\end{teobox}

A intuição do teorema é direta e vale mais que o enunciado. Para ter $R_n$
pequeno, é preciso \emph{parar} de amostrar os braços que parecem ruins — e é
exatamente amostrá-los que traria certeza sobre qual é o melhor. Os dois objetivos
puxam em direções opostas.

\subsection{O problema metanível}

A pergunta correta não é ``qual ação simular para ganhar mais recompensa
simulada'', e sim:

\begin{teobox}{Formulação metanível}
Qual é a próxima \emph{computação} a executar, dado que ela custa tempo e que seu
único benefício é a chance de \emph{mudar} a decisão da raiz?
\end{teobox}

Hay, Russell, Tolpin e Shimony (2012) mostram que isto é, literalmente, um MDP:
os estados são estados de crença sobre os valores das ações, as ações são
computações, e a recompensa é o valor da decisão final menos o custo do tempo
gasto. A consequência prática mais útil é um critério de valor da informação:

\begin{center}
\emph{Uma computação só vale a pena se puder mudar a jogada escolhida.}
\end{center}

Refinar a estimativa do braço que já está em último lugar é desperdício, por mais
incerta que essa estimativa seja — a incerteza não é o critério, a possibilidade
de reversão do ranking é. Algoritmos desenhados para esse objetivo — BRUE, SHOT, e
variantes de \emph{sequential halving} em árvores — superam o UCT em vários
cenários de orçamento fixo.

\begin{ideiabox}
Curiosamente, o PUCT do AlphaZero se afasta do UCT \emph{na direção certa}, ainda
que por acidente histórico. O decaimento $\Nvis(s,a)^{-1}$ é mais agressivo que o
$\Nvis(s,a)^{-1/2}$ do UCT, e o crescimento $\sqrt{\Nvis(s)}$ é muito mais forte
que $\sqrt{\log \Nvis(s)}$. O efeito líquido é distribuir visitas de forma mais
parecida com uma \emph{política} do que com uma cota de confiança — que é
precisamente o que a leitura de Grill et al.\ formaliza.
\end{ideiabox}

\subsection{A parede da profundidade: um experimento}

O argumento teórico fica mais convincente com um exemplo executável. Considere a
\textbf{árvore-armadilha}, no espírito das construções de Coquelin e Munos (2007):

\begin{algbox}{Árvore-armadilha de profundidade $D$}
Em cada profundidade $d \in \{0,\dots,D-1\}$ há duas ações:
\begin{itemize}
  \item \texttt{parar}: termina imediatamente com valor $0{,}9 - 0{,}4\,d/D$;
  \item \texttt{seguir}: vai para a profundidade $d+1$.
\end{itemize}
Chegar a $d = D$ vale $1{,}0$. A política ótima é seguir sempre, com valor
$1{,}0$; a tentação está no começo, e a recompensa está no fim. Rollouts a partir
de um nó folha escolhem \texttt{parar} ou \texttt{seguir} com probabilidade
$1/2$, de modo que alcançam a folha profunda com probabilidade $2^{-(D-d)}$.
\end{algbox}

Medindo a probabilidade de a raiz recomendar a ação ótima, com MCTS/UCT clássica e
escolhendo o melhor $c$ para cada célula:

\begin{center}
\begin{tabular}{ccccc}
\toprule
$D$ & $K=200$ & $K=1\,000$ & $K=5\,000$ & $K=20\,000$ \\
\midrule
4  & 0,85 & 1,00 & 1,00 & 1,00 \\
6  & 0,00 & 0,20 & 1,00 & 1,00 \\
8  & 0,00 & 0,00 & 0,00 & 0,00 \\
10 & 0,00 & 0,00 & 0,00 & 0,00 \\
\bottomrule
\end{tabular}
\end{center}

Com $D = 8$ nem $50\,000$ simulações resolvem. Duas causas se somam.

\textbf{(i) A rollout aleatória quase nunca vê o valor bom.} A probabilidade de
alcançar a folha profunda decai como $2^{-D}$. O que a busca observa, na esmagadora
maioria das simulações, são os valores de \texttt{parar}.

\textbf{(ii) O UCT retropropaga a média, não o máximo.} Como observado na
Seção~1.5, a média é tomada sobre uma política não estacionária. Ela fica
contaminada por exploração antiga, e o atraso de cada nível se compõe ao subir.
Coquelin e Munos exibem famílias de árvores em que o UCT precisa de um número de
amostras \emph{exponencial em uma exponencial} na profundidade. A garantia de
consistência assintótica é verdadeira e inútil.

\subsection{O ajuste de $c$ não é secundário}

Mesma árvore com $D = 6$ e orçamento fixo de $2\,000$ simulações, variando apenas
a constante de exploração:

\begin{center}
\begin{tabular}{lccccccc}
\toprule
$c$ & 0,1 & 0,2 & 0,3 & 0,5 & \textbf{0,7} & 1,0 & 1,4 \\
\midrule
acerto & 0,000 & 0,000 & 0,125 & 0,270 & \textbf{0,650} & 0,000 & 0,000 \\
\bottomrule
\end{tabular}
\end{center}

A curva não é monotônica nem suave. Pouca exploração nunca sai da armadilha rasa;
muita exploração contamina de tal modo as médias retropropagadas que a raiz nunca
enxerga a vantagem do caminho profundo. Note também que $c$ precisa ser calibrado
à \emph{escala dos valores}: aqui os retornos vivem em $[0{,}5;\,1{,}0]$, e um
$c = 1{,}4$ que seria razoável para recompensas em $[0,10]$ é aqui destrutivo.

Isto explica, em retrospecto, os 94 e 103 pontos de Elo obtidos no AlphaGo apenas
com otimização bayesiana dos hiperparâmetros da busca.

\subsection{Então por que a MCTS funciona tão bem na prática?}

Quatro razões, em ordem de importância.

\textbf{1. A rollout foi eliminada.} A causa (i) da parede da profundidade
desaparece quando a folha é avaliada por $\redeval$ em vez de por uma partida
aleatória. O bom valor não precisa mais ser \emph{encontrado} por sorte: ele é
\emph{previsto}.

\textbf{2. O prior poda a largura.} Das $\approx 250$ jogadas legais no Go, o PUCT
examina seriamente uma dúzia. O fator de ramificação efetivo despenca.

\textbf{3. As árvores reais não são adversárias.} Nas armadilhas construídas por
Coquelin e Munos o valor está escondido atrás de um caminho estreito e único; em
Go e xadrez, posições boas costumam ter vizinhas boas. Existe gradiente a seguir.

\textbf{4. O laço de treino conserta a heurística.} Onde $\redeval$ erra
sistematicamente, o auto-jogo gera dados exatamente ali, porque a busca continua
levando o agente àquelas posições.

\begin{teobox}{A leitura para levar}
A MCTS moderna funciona \emph{apesar} de o UCT ser questionável, e não por causa
dele. O peso saiu da regra de seleção e foi para a qualidade do avaliador de
folha. É por isso que ``mais simulações'' rende cada vez menos, enquanto ``rede
melhor'' continua rendendo. Consequência prática: ao diagnosticar um planejador
que decide mal, suspeite primeiro do avaliador de folha e do prior, e só depois da
constante de exploração ou do orçamento.
\end{teobox}

\begin{discbox}
O critério metanível diz que uma computação só vale a pena se puder mudar a
decisão. Aplicado ao AlphaZero, ele sugeriria interromper a busca assim que a
jogada mais visitada se tornasse inalcançável pelas demais dentro do orçamento
restante — uma parada antecipada barata e correta. Por que isso raramente é feito?
E qual seria o análogo dessa ideia em um modelo de linguagem que ``pensa'' por
mais tempo antes de responder?
\end{discbox}

%% ============================================================
\section{Formulário}

\begin{center}
\begin{tabular}{@{}p{0.30\linewidth} p{0.66\linewidth}@{}}
\toprule
\textbf{Objeto} & \textbf{Expressão} \\
\midrule
Busca MC simples &
  $\widehat{Q}(s_t,a) = \frac{1}{K}\sum_k G^k_t$;\quad
  $a_t = \argmax_a \widehat{Q}(s_t,a)$ \\[0.35em]
Tamanho do expectimax &
  $(\abs{\gS}\abs{\gA})^{H}$ \\[0.35em]
Seleção UCT &
  $\argmax_a \bigl[\overline{G}(s,a) + c\sqrt{\log \Nvis(s)/\Nvis(s,a)}\bigr]$ \\[0.35em]
Seleção PUCT &
  $\argmax_a \bigl[Q(s,a) + \cpuct P(s,a)\sqrt{\Nvis(s)}/(1+\Nvis(s,a))\bigr]$ \\[0.35em]
Retropropagação &
  $\Nvis \mathrel{+}= 1$;\ \ $\Wtot \mathrel{+}= \pm v$;\ \ $Q = \Wtot/\Nvis$ \\[0.35em]
Ruído na raiz &
  $P(s_0,a) = (1-\varepsilon)\redepol(s_0)_a + \varepsilon \eta_a$,\ \
  $\eta\sim\mathrm{Dir}(\alpha)$,\ $\varepsilon = 0{,}25$ \\[0.35em]
Política da raiz &
  $\raizpol(a\mid s_0) = \Nvis(s_0,a)^{1/\tau} / \sum_b \Nvis(s_0,b)^{1/\tau}$ \\[0.35em]
Perda do AlphaZero &
  $\ell(\theta) = (z-\redeval(s))^2 - \raizpol^{\top}\log\redepol(s) + c\norm{\theta}^2$ \\[0.35em]
Forma fechada de Grill &
  $\bar{\raizpol}_a = \lambda_{\Nvis}\,p_{\theta,a}/(\alpha - q_a)$,\ \
  $\lambda_{\Nvis} = \cpuct\sqrt{\Nvis}/(\Nvis+\abs{\gA})$ \\[0.35em]
Arrependimento simples &
  $r_n = \mu^{*} - \mu_{\widehat{A}_n}$ \quad (contra $R_n = \sum_t (\mu^{*}-\mu_{A_t})$) \\[0.35em]
Escala Elo &
  $p(\text{vitória}) = \bigl(1 + 10^{-\Delta/400}\bigr)^{-1}$;\ \ $\Delta = 200 \Rightarrow p \approx 0{,}76$ \\
\bottomrule
\end{tabular}
\end{center}

\begin{center}
\begin{tabular}{@{}lll@{}}
\toprule
\textbf{Sistema} & \textbf{Avaliação da folha} & \textbf{Modelo} \\
\midrule
MCTS clássica / UCT & rollout aleatória & dado (simulador) \\
AlphaGo (2016)      & $\lambda\,$rollout $+\,(1-\lambda)\,\redeval$ & dado \\
AlphaGo Zero (2017) & $\redeval$ apenas & dado \\
AlphaZero (2018)    & $\redeval$ apenas & dado, três jogos \\
MuZero (2020)       & $\redeval$ apenas & \textbf{aprendido}, latente \\
\bottomrule
\end{tabular}
\end{center}

%% ============================================================
\section{Exercícios}

\begin{exbox}{L14E1 — PUCT à mão \normalfont(com resposta)}
Raiz com três ações, priors $P = (0{,}5;\ 0{,}3;\ 0{,}2)$, valores das folhas
$\redeval = (0{,}20;\ 0{,}60;\ 0{,}10)$ e $\cpuct = 1{,}0$. Use a convenção
$\sqrt{\Nvis(s)} \to 1$ na primeira simulação. Execute oito simulações à mão,
preenchendo $Q$, $U$ e $Q+U$ a cada passo, e reporte $\Nvis$ final e
$\raizpol$ com $\tau = 1$.
\end{exbox}

\begin{provabox}
\textbf{Resposta.} A sequência de ações escolhidas é
$a_1, a_1, a_1, a_2, a_2, a_2, a_2, a_2$.

\begin{center}
\small
\begin{tabular}{c ccc c ccc c c}
\toprule
& \multicolumn{3}{c}{$Q$} && \multicolumn{3}{c}{$U$} && \\
\cmidrule(lr){2-4}\cmidrule(lr){6-8}
$k$ & $a_1$ & $a_2$ & $a_3$ && $a_1$ & $a_2$ & $a_3$ && escolhe \\
\midrule
1 & 0,000 & 0,000 & 0,000 && 0,500 & 0,300 & 0,200 && $a_1$ \\
2 & 0,200 & 0,000 & 0,000 && 0,250 & 0,300 & 0,200 && $a_1$ \\
3 & 0,200 & 0,000 & 0,000 && 0,236 & 0,424 & 0,283 && $a_1$ \\
4 & 0,200 & 0,000 & 0,000 && 0,217 & 0,520 & 0,346 && $a_2$ \\
5 & 0,200 & 0,600 & 0,000 && 0,250 & 0,300 & 0,400 && $a_2$ \\
6 & 0,200 & 0,600 & 0,000 && 0,280 & 0,224 & 0,447 && $a_2$ \\
7 & 0,200 & 0,600 & 0,000 && 0,306 & 0,184 & 0,490 && $a_2$ \\
8 & 0,200 & 0,600 & 0,000 && 0,331 & 0,159 & 0,529 && $a_2$ \\
\bottomrule
\end{tabular}
\end{center}

Final: $\Nvis = (3,\ 5,\ 0)$ e $\raizpol = (0{,}375;\ 0{,}625;\ 0)$.

\textbf{Observe} que $a_1$ é escolhida três vezes seguidas apesar de já ter
devolvido $0{,}20$ na primeira: com $\cpuct$ menor e prior mais concentrado que no
exemplo da Seção~\ref{sec:traco}, a busca demora mais a migrar. E $a_3$ nunca é
visitada, embora seu $U$ seja o maior dos três a partir de $k=6$ — o que falta é
orçamento, não incentivo.
\end{provabox}

\begin{exbox}{L14E2 — Os dois regimes da forma fechada}
Usando $\bar{\raizpol}_a = \lambda\,p_a/(\alpha - q_a)$ com $\alpha$ resolvendo
$\sum_a \lambda p_a/(\alpha - q_a) = 1$:
\begin{enumerate}
  \item Mostre que $\lambda \to \infty$ implica $\bar{\raizpol} \to \bm{p}$.
  \item Mostre que $\lambda \to 0^{+}$ implica $\bar{\raizpol} \to$ indicadora do
        $\argmax$ de $\bm{q}$ (suponha o máximo único).
  \item Com $p = (0{,}6;\ 0{,}3;\ 0{,}1)$ e $q = (0{,}10;\ 0{,}55;\ -0{,}40)$,
        calcule numericamente $\alpha$ e $\bar{\raizpol}$ para $\lambda = 0{,}0374$
        e confira contra a tabela da Seção~2.8.
\end{enumerate}
\end{exbox}

\begin{exbox}{L14E3 — Temperatura e entropia}
Com $\Nvis = (5,\ 11,\ 0)$:
\begin{enumerate}
  \item Calcule $\raizpol$ e $H[\raizpol]$ para $\tau \in \{1;\ 0{,}5;\ 0{,}25;\ 0{,}1\}$.
  \item Mostre que $\lim_{\tau\to 0^{+}} \raizpol$ é a indicadora do $\argmax_a \Nvis$
        e que $\lim_{\tau\to\infty} \raizpol$ é uniforme sobre as ações
        \emph{visitadas}.
  \item O AlphaZero usa $\tau = 1$ nas 30 primeiras jogadas. Argumente por que
        usar $\tau = 1$ na partida \emph{inteira} degradaria simultaneamente a
        força de jogo e a qualidade do alvo de valor.
\end{enumerate}
\end{exbox}

\begin{exbox}{L14E4 — Densidade de informação no auto-jogo}
Modele o resultado de uma partida como Bernoulli$(p)$, com $p$ a probabilidade de
vitória do agente que está treinando.
\begin{enumerate}
  \item Calcule $H(p)$ em bits para $p \in \{0{,}5;\ 0{,}75;\ 0{,}9;\ 0{,}99\}$ e
        o número de partidas necessárias para acumular 1 bit.
  \item Suponha um adversário fixo contra o qual o agente vence 99\% das vezes.
        Quantas vezes mais partidas são necessárias, em relação ao auto-jogo, para
        obter a mesma informação sobre o parâmetro?
  \item O argumento acima trata só do alvo de \emph{valor}. Explique por que o
        alvo de \emph{política} não sofre a mesma degradação — e por que, ainda
        assim, um adversário desbalanceado prejudica o treino.
\end{enumerate}
\end{exbox}

\begin{exbox}{L14E5 — Elo e o valor da busca}
Use $p = (1 + 10^{-\Delta/400})^{-1}$.
\begin{enumerate}
  \item Confirme que $\Delta = 200$ dá $p \approx 0{,}76$.
  \item A rede crua do AlphaGo Zero tem Elo $3\,055$ e o sistema completo,
        $5\,185$. Calcule $p$ e o número esperado de partidas até a rede crua
        vencer uma.
  \item A escala Elo é logarítmica na razão de chances. Discuta em que sentido
        dizer que ``a busca vale $2\,130$ Elo'' é mais informativo — e em que
        sentido é menos — do que dizer que ela vence $99{,}9995\%$ das partidas.
\end{enumerate}
\end{exbox}

\begin{exbox}{L14E6 — A árvore-armadilha \normalfont(implementação)}
Implemente a árvore-armadilha do algoritmo da Seção~6.3 e a MCTS/UCT com rollout
aleatória.
\begin{enumerate}
  \item Reproduza a tabela de acerto na raiz para $D \in \{4,6,8\}$ e
        $K \in \{200;\ 1\,000;\ 5\,000\}$, varrendo $c \in \{0{,}3;\ 0{,}5;\ 0{,}7;\ 1{,}0\}$.
  \item Substitua a rollout por um avaliador $\widehat{v}(d) = V^{*}(d) + \epsilon$
        com $\epsilon \sim \mathcal{N}(0,\sigma^2)$. Trace a curva de acerto
        \emph{versus} $\sigma$ para $D = 8$ e $K = 1\,000$. A partir de que $\sigma$
        o avaliador deixa de ajudar?
  \item Instrumente o código para registrar, a cada 100 simulações, o valor
        $Q(\text{raiz},\texttt{seguir})$. Verifique empiricamente que ele converge
        para $1{,}0$ \emph{por baixo}, e relacione isso com a discussão sobre
        média \emph{versus} máximo da Seção~1.5.
\end{enumerate}
\end{exbox}

\begin{exbox}{L14E7 — Min-max implícito}
No AlphaZero, $\redeval(s)$ é sempre o valor do ponto de vista do jogador da vez, e
a retropropagação alterna o sinal a cada nível.
\begin{enumerate}
  \item Escreva a recorrência de $Q(s,a)$ resultante e mostre que, no limite de
        muitas simulações, a árvore implementa uma avaliação min-max.
  \item Onde exatamente essa construção usa a hipótese de \emph{soma zero}?
        Descreva o que quebra em um jogo de dois jogadores com recompensas
        arbitrárias.
  \item Como você adaptaria o esquema para três jogadores? (Dica: pense no que
        substitui o par $(\max,\min)$.)
\end{enumerate}
\end{exbox}

\begin{exbox}{L14E8 — Quando \emph{não} usar MCTS}
Para cada cenário, diga se a MCTS é uma boa escolha e justifique em uma frase.
\begin{enumerate}
  \item Controle de um pêndulo invertido, estado contínuo em $\R^4$, ação
        contínua, horizonte de 500 passos, simulador disponível.
  \item Escalonamento de tarefas com 8 máquinas e 20 tarefas, modelo exato
        conhecido, horizonte 20.
  \item Recomendação de conteúdo com milhões de itens, sem simulador confiável do
        comportamento do usuário.
  \item Planejamento de rotas em um grafo de $10^{7}$ nós com custos conhecidos e
        determinísticos.
\end{enumerate}
\end{exbox}

\begin{exbox}{L14E9 — Custo de decisão \emph{versus} custo de treino}
O AlphaGo Zero usa $1\,600$ simulações por jogada, a $\approx 0{,}4$\,s por busca,
e gerou $4{,}9$ milhões de partidas em três dias. Supondo 200 jogadas por partida:
\begin{enumerate}
  \item Estime o número total de avaliações da rede na fase de auto-jogo.
  \item Compare com o número de passos de gradiente ($\approx 700\,000$ no
        AlphaZero, com lotes de $4\,096$). Qual fase domina o custo?
  \item Que consequência isso tem para quem quer reproduzir o experimento com
        orçamento acadêmico? Que parâmetro você reduziria primeiro, e o que se
        perde?
\end{enumerate}
\end{exbox}

%% ============================================================
\section{Autoavaliação}

Marque o que você consegue fazer sem consultar o material.

\begin{enumerate}
  \item Explicar por que jogar Go \emph{não} é um problema de modelo desconhecido,
        e o que de fato torna o Go difícil.
  \item Escrever a regra PUCT de memória e dizer o papel de cada um dos quatro
        fatores do termo $U$.
  \item Listar três diferenças entre UCT e PUCT, e dizer qual delas explica por
        que o AlphaZero precisa de ruído de Dirichlet.
  \item Explicar por que trocar a rollout pela rede de valor é a mesma troca de
        Monte Carlo por TD, incluindo o efeito sobre viés e variância.
  \item Justificar por que a política da raiz usa contagens de visita e não
        $\argmax_a Q$.
  \item Derivar a forma fechada $y_a = \lambda p_a/(\alpha - q_a)$ e identificar os
        dois regimes-limite.
  \item Explicar em que sentido a MCTS é um operador de melhoria de política, e
        por que isso transforma o laço de treino em iteração de política.
  \item Argumentar, com o cálculo de entropia, por que o auto-jogo mantém a
        recompensa densa em informação.
  \item Distinguir arrependimento simples de cumulativo e explicar por que o UCT
        otimiza o segundo dentro de um problema que pede o primeiro.
  \item Descrever a parede da profundidade, suas duas causas, e por que a rede de
        valor elimina uma delas.
\end{enumerate}

%% ============================================================
\section{Leituras}

\textbf{Fontes primárias}
\begin{itemize}
  \item Silver, D.\ et al. \emph{Mastering the game of Go with deep neural networks
        and tree search}. Nature 529, 484--489, 2016.
  \item Silver, D.\ et al. \emph{Mastering the game of Go without human knowledge}.
        Nature 550, 354--359, 2017.
  \item Silver, D.\ et al. \emph{A general reinforcement learning algorithm that
        masters chess, shogi and Go through self-play}. Science 362, 1140--1144, 2018.
  \item Schrittwieser, J.\ et al. \emph{Mastering Atari, Go, chess and shogi by
        planning with a learned model}. Nature 588, 604--609, 2020.
  \item Fawzi, A.\ et al. \emph{Discovering faster matrix multiplication algorithms
        with reinforcement learning}. Nature 610, 47--53, 2022.
  \item Mankowitz, D.\ et al. \emph{Faster sorting algorithms discovered using deep
        reinforcement learning}. Nature 618, 257--263, 2023.
\end{itemize}

\textbf{Fundamentos, análise e crítica}
\begin{itemize}
  \item Kocsis, L.\ e Szepesvári, C. \emph{Bandit based Monte-Carlo planning}.
        ECML, 2006. \hfill (UCT)
  \item Coquelin, P.-A.\ e Munos, R. \emph{Bandit algorithms for tree search}.
        UAI, 2007. \hfill (patologias do UCT)
  \item Bubeck, S., Munos, R.\ e Stoltz, G. \emph{Pure exploration in multi-armed
        bandits problems}. ALT, 2009.
  \item Hay, N., Russell, S., Tolpin, D.\ e Shimony, S. \emph{Selecting
        computations: theory and applications}. UAI, 2012.
  \item Browne, C.\ et al. \emph{A survey of Monte Carlo tree search methods}.
        IEEE TCIAIG 4(1), 1--43, 2012.
  \item Anthony, T., Tian, Z.\ e Barber, D. \emph{Thinking fast and slow with deep
        learning and tree search}. NeurIPS, 2017.
  \item Chen, Y.\ et al. \emph{Bayesian optimization in AlphaGo}.
        arXiv:1812.06855, 2018.
  \item Grill, J.-B.\ et al. \emph{Monte-Carlo tree search as regularized policy
        optimization}. ICML, 2020.
\end{itemize}

\vfill
\begin{center}
\footnotesize\color{rlCinza}
Material adaptado de \emph{CS234: Reinforcement Learning}, Emma Brunskill,
Stanford University.\\
Universidade Federal de Uberlândia \textbullet\ Prof.\ Marcelo Keese Albertini
\end{center}

\end{document}
